\documentclass[10pt,a4paper]{amsart}
\usepackage[margin=19mm]{geometry}
\usepackage{amsmath,amssymb,mathtools}
\usepackage[scr=rsfs,cal=euler]{mathalfa}
\usepackage{xfrac}
\usepackage[unicode,hidelinks,bookmarks=false]{hyperref}
\newcommand{\ie}{i.e.\ }
\newcommand{\cf}{cf.\ }
\newcommand{\resp}{respectively\ }
\newcommand{\Subsection}{\subsection}
\providecommand{\nobreakdash}{\nobreak\hbox{-}}
\setlength{\emergencystretch}{2em}
\multlinegap0pt
\usepackage{tikz}
\usetikzlibrary{cd}
\usepackage{amssymb}
\usepackage[shortlabels]{enumitem}
\usepackage{xcolor}
\usepackage{tikz}
\newtheorem{theorem}{Theorem}
\newtheorem{lemma}{Lemma}
\newcommand{\F}{\mathbb{F}}
\newcommand{\Fbar}{{\overline{\F}}}
\DeclareMathOperator{\Frob}{Frob}
\newcommand{\PP}{\mathbb{P}}
\title{Exponents of Jacobians and relative class groups}
\author{Borys Kadets, Daniel Keliher}
\date{Source: arXiv:2410.11962v2 (CC BY 4.0)}
\begin{document}
\maketitle

\noindent\emph{Source and licence.} Exponents of Jacobians and relative class groups. Version arXiv:2410.11962v2 (CC BY 4.0), distributed by arXiv under the Creative Commons Attribution 4.0 International licence, http://creativecommons.org/licenses/by/4.0/, as recorded in the arXiv licence metadata for this exact version and shown on the abstract page https://arxiv.org/abs/2410.11962v2. Full licence notice: \texttt{licenses/arxiv-2410.11962-CC-BY-4.0.txt}.\par\smallskip

\noindent\textit{Editorial scope.} Counted unit: the article's lemma lem:non-fibralcount (source0.tex lines 347--351). The article's complete original proof of it is delivered with it (source0.tex lines 353--364, one \texttt{proof} environment of 12 lines). No mathematics is written, repaired or re-derived: the statement and the proof are the article's own lines, in the article's own notation, and no retained line is substituted or re-flowed.  Retained uncounted as the source-local context the proof needs: lines 332--336.  Nothing was omitted from any retained span, so the package carries no omission record.  No retained line contains a citation, so the wrapper carries no bibliography and no bibliography entry.  The retained lines contain the article's own diagram code, which the wrapper typesets with the standard tikz packages; no image file is delivered and the package declares no asset dependency.  Editorial formatting only, in addition: an AMS \texttt{amsart} preamble that restates the article's own declarations and macros used by the retained lines, namely the environments \texttt{theorem}, \texttt{lemma} on separate counters, together with the standard packages those lines need.  The retained environments print in this document as Theorem 1, Lemma 1 in order of appearance, whereas the article prints the same items with the numbering of its own full document.  The complete native source of the article as distributed in the arXiv e-print for this exact version is bundled unchanged as \texttt{sources/166895/source0.tex}.
\begin{theorem}[Riemann--Hurwitz inequality]\label{thm:RH}
    Suppose $\phi: X \to Y$ is a separable covering of curves. For a point $P \in |X|$ let $e(P)$ denote the ramification index of $\phi$ at $P$. Then the genera $g_X, g_Y$ of $X$ and $Y$ satisfy 
    \[2g_X-2 \geqslant \deg \phi(2g_Y-2)+\sum_{P \in |X|}\deg P(e(P)-1).\]
    In particular, the number of branch points in $X(\Fbar_q)$ is bounded from above by $(2g_X-2)-\deg \phi (2g_Y-2).$
\end{theorem}

\begin{lemma}\label{lem:non-fibralcount}
    If $X/\F_q$ is a curve of genus $g>0$, $f: X \to \PP^1$ is a rational function, and $k$ is a prime number, then there are at least 
    \[\left\lceil\frac{q^k}{k} - 2g\frac{q^{k/2}+1}{k} - \deg f \frac{q+3}{k}\right\rceil\] 
    degree $k$ points $P \in |X|$ which, when viewed as degree $k$ divisors, are non-fibral under $f$. 
\end{lemma}

\begin{proof}
    Since $g$ is greater than $0$, $f$ is not a power of the Frobenius. Thus $f$ factors as follows
    \begin{center}
        \begin{tikzcd}
        X \arrow[r, "\mathrm{Frob}_{p^n}"] \arrow[bend right=30]{rr}{f} & X^{(p^n)} \arrow[r, "f'"] & \mathbb{P}^1,
        \end{tikzcd}
    \end{center}
    where $f':X^{(p^n)} \to \PP^1$ is a separable, nonconstant map from the Frobenius twist. Since $\Frob_{p^n}$ is a bijection between degree $k$ points on $X$ and $X^{(p^n)}$, it is enough to prove the result for the separable morphism $f'$.

    By the Weil conjectures, there are at least $q^k-2gq^{k/2}+1$ points in $X(\F_{q^k})$. At most $(\deg f') (q+1)$ among these points are sent to $\PP^1(\F_q)$ by $f'$. Moreover, since $f'$ is separable, by the Riemann--Hurwitz formula (Theorem \ref{thm:RH}), at most $2\deg f' + 2g -2$ points on $X$ are ramified. Thus we can choose $N$ points among $X(\F_{q^k})$, whose image under $f'$ does not belong to $\PP^1(\F_q)$, and which do not intersect the ramification locus of $f'$ for \[N=(q^k-2gq^{k/2}+1) - \deg f' (q+1) - (2\deg f' + 2g - 2).\]
    Note that, since $k$ is prime, any point in $X(\F_{q^k})\setminus X(\F_q)$, in particular those among the $N$ chosen, come from a degree $k$ point in $|X|$. Passing to Galois orbits, we have at least $\lceil N/k \rceil$ non-fibral degree $k$ points in $|X|$.
\end{proof}

\end{document}
